\appendix
\section{Second-Moment and Transport Identities}
For normalized point deposition, the new memory is a mixture of the old
measure with weight $q$ and a point measure with weight $\alpha$. The
variance identity for this mixture is
\begin{align}
Q_{n+1}&=q Q_n+q|c_n-c_{n+1}|^2
+\alpha|x_{n+1}-c_{n+1}|^2\\
&=qQ_n+q\alpha|x_{n+1}-c_n|^2.
\end{align}
This proves Eq.~(\ref{eq:cloud_recursion}) without a Gaussian assumption.
Its stationary expectation and $r_{n+1}=q(x_{n+1}-c_n)$ give
Eq.~(\ref{eq:cloud_radius}). The position-relative variance follows from
the innovation independence in Eq.~(\ref{eq:relative_mode}). For a centered
deposition profile of nonzero variance, the cloud recursion acquires its
deposition variance term; the point-deposition result must not be applied
unchanged to that model.

For finite history, let $h_n$ be the visible-coordinate response to a unit
innovation, with zero prehistory. The centroid response is
$b_n=\sum_jw_jh_{n-j}$. The two impulse energies per coordinate are
\begin{align}
E_r&=\sum_{n\geq0}(h_n-b_n)^2,\\
E_Q&=\sum_{n\geq0}\sum_jw_j(h_{n-j}-b_n)^2.
\end{align}
Multiplying by $d\varepsilon^2$ yields the second moments. The position
response itself approaches a nonzero constant, but both centered energies
converge. Parseval's identity gives Eqs.~(\ref{eq:finite_relative}) and
(\ref{eq:finite_cloud}); for the latter the weighted sum of squared filter
differences simplifies to $1-|B_H|^2$.

For $H>1$ and $0<g<1$, the noiseless position-history companion map has
nonnegative top-row coefficients $(1-g+gw_0,gw_1,\ldots,gw_{H-1})$ summing
to one and shift rows below it. The matrix is irreducible and aperiodic.
Its eigenvalue one is the translation mode, and all remaining eigenvalues
have modulus below one. At $g=0$, the visible process is a random walk;
its finite-window lag differences still have stationary second moments.
This justifies the relative linear reference for the tested gain cells;
it is not a stability theorem for the nonlinear Gaussian map.

Expanding the filter near zero frequency gives
$B_H(e^{i\omega})=1-i\omega\bar j_H+O(\omega^2)$ and
$F_H(e^{i\omega})=i\omega(1+g\bar j_H)+O(\omega^2)$.
Thus the large-scale innovation variance is
$\varepsilon^2/(1+g\bar j_H)^2$, proving Eq.~(\ref{eq:diffusion}), with
\begin{equation}
\bar j_H=\frac{q}{\alpha}-\frac{Hq^H}{1-q^H}.
\end{equation}
The untruncated limit has $\bar j=q/\alpha$. The discrete update count
remains a model parameter; no physical-time calibration is inferred.

\section{Reproduction and Measurement Limits}
The finite reference is evaluated by two implementations: a causal FIFO
recurrence with a single unit innovation and a midpoint frequency integral
with 131,072 points on $[0,\pi]$. The reference tests include $H=1$, the
zero-gain random-walk covariance obtained from pairwise lag distances, the
untruncated limit, and a second $(\alpha,H,g)$ cell distinct from the
reported long runs. Their purpose is to test the reference calculation,
not to provide new nonlinear simulation evidence.

The terminal block is selected from the stored integer timestamps before
any radius ratio is computed. It is the maximal regular suffix, which is
10,001 measurements at unit stride for every reported case. Source-file
hashes and the extracted-array archive hash are recorded; a rebuild checks
the archive hash. Squared radii are averaged first, then square-rooted;
the seed ratios are aggregated last. No instantaneous-radius median is
substituted into the RMS comparison.

The two halves of each terminal block are also reported in the
machine-readable output. They expose finite-window variability without
selecting a favorable subwindow. Since these measurements are correlated,
neither the number of stored points nor the preceding $N$ updates supplies
an independent sample count. The five-seed IQR is the only uncertainty
summary used in the main figures. No statistical equivalence margin or
confidence-level claim is attached to the observed discrepancies.
